% Software Requirements Specification following the Meyer PEGS template
% (Goals, Environment, System, Project) from Bertrand Meyer's "Handbook of
% Requirements and Business Analysis" (2022).
%
% The LaTeX structure is adapted from the Meyer-template SRS written by the
% AlgoCatan team, shared with us by our TA:
% https://github.com/AlgoCatan/RLCatan/blob/main/docs/SRS/SRS.tex
% Their project content has been removed; only the structure, numbering
% mechanics and table layouts are kept. Deviations from Meyer's original
% template are listed in "Changes from Original SRS Meyer Template".
%
% Template guidance is written with \plt{...} (cyan). Delete each note once
% the section it describes has been written.

\documentclass[12pt]{article}

\usepackage{tabularx}
\usepackage{booktabs}
\usepackage{array}
\usepackage{amsfonts}
\usepackage{amsmath}
\usepackage{amstext}
\usepackage{float}
\usepackage{graphicx}
\usepackage{zed-csp}     % Z notation, for S.9 Formalizations (optional)
\usepackage{placeins}
\usepackage{pdflscape}
\usepackage{caption}
\usepackage[round]{natbib}

\input{../Comments.text}
\input{../Common.text}

% ---------------------------------------------------------------------------
% PEGS chapter macro. Meyer's template numbers subsections as G.1, E.2, S.3,
% P.4 etc. Each chapter is an unnumbered section with a manual TOC entry, and
% the subsection counter is reset with the chapter letter as its prefix.
% ---------------------------------------------------------------------------
\newcommand{\pegschapter}[2]{%
  \newpage
  \phantomsection
  \section*{(#1) #2}\label{sec:srs-#1}
  \addcontentsline{toc}{section}{(#1) #2}
  \renewcommand{\thesubsection}{#1.\arabic{subsection}}
  \renewcommand{\theHsubsection}{#1.\arabic{subsection}}
  \setcounter{subsection}{0}%
}

% Unnumbered section that still appears in the table of contents.
\newcommand{\tocsection}[1]{%
  \phantomsection
  \section*{#1}
  \addcontentsline{toc}{section}{#1}%
}

\title{Software Requirements Specification for \progname: subtitle describing software}
\author{\authname}
\date{\today}

\begin{document}

\maketitle
\thispagestyle{empty}

~\newpage

\pagenumbering{roman}

\begin{table}[hp]
\caption{Revision History} \label{TblRevisionHistory}
\begin{tabularx}{\textwidth}{llX}
\toprule
\textbf{Date} & \textbf{Developer(s)} & \textbf{Change}\\
\midrule
2026-10-09 & Colin Chambachan & Switch from the Volere template to the Meyer template\\
Date2 & Name(s) & Description of changes\\
... & ... & ...\\
\bottomrule
\end{tabularx}
\end{table}

~\newpage

\tableofcontents

\listoftables

\listoffigures

~\newpage

\pagenumbering{arabic}

% ===========================================================================
\tocsection{Changes from Original SRS Meyer Template}

\plt{Meyer's template is a book chapter, not a course template, so every team
adapts it. Record each structural change you make and why, so a reader who
knows the original can follow this document. The rows below describe the
adaptations already present in this file; add to them as the document evolves.}

\begin{table}[H]
\caption{Changes from Template} \label{ChangesFromTemplate}
\begin{tabularx}{\textwidth}{|p{4.5cm}|>{\raggedright\arraybackslash}X|}
\hline
\textbf{Change} & \textbf{Reason} \\
\hline
Glossary (E.1) moved to the front of the document & Terms and acronyms must be defined before first use, and information should live in one place. E.1 now points here. \\
\hline
Requirement tables added to each PEGS chapter & Meyer's template describes requirements in prose. Labelled tables of functional (FR) and non-functional (NFR) requirements give each requirement an ID that can be prioritized, traced and verified. \\
\hline
Added G.8 User Profiles & Meyer lists stakeholders (G.7) but has no place for the distinct characteristics of each user group. \\
\hline
Added E.7 Standards and Legal/Regulatory Factors & Meyer's Environment chapter has no explicit home for standards, licensing or legal constraints. \\
\hline
Added S.4 Normal Operation and S.5 Undesired Event Handling & Separates expected behaviour from error handling instead of mixing both into S.2 Functionality. \\
\hline
Added S.9 Formalizations & Provides a place for mathematical or formal specification of the parts of the system that benefit from it. \\
\hline
Added Safety and Security Requirements & Lets the requirements from the Hazard Analysis be ported into this document. \\
\hline
Added Requirement Dependencies & Shows traceability between requirements visually. \\
\hline
Added References & Meyer's template has no references section. \\
\hline
Added Appendices (Generative AI Use Disclosure, Reflection) & Required by the course. \\
\hline
\end{tabularx}
\end{table}

% ===========================================================================
\tocsection{Glossary}
\label{sec:glossary}

\plt{Define every domain term and acronym used in this document, in
alphabetical order. Define a term once, here, and use it consistently
everywhere else. Where a term has a good external definition, link to it and
cite the source.}

\plt{Tip from the source template: define a \texttt{\textbackslash hypertarget}
for each glossary entry and a macro such as \texttt{\textbackslash newcommand\{\textbackslash Term\}\{\textbackslash hyperlink\{glossary-term\}\{Term\}\}}
so each use of the term in the body links back to its definition.}

\begin{itemize}
    \item \textbf{Term} -- Definition.
    \item \textbf{Acronym (Expansion)} -- Definition.
\end{itemize}

% ===========================================================================
\pegschapter{G}{Goals}

\plt{The Goals chapter is written for decision makers. It explains why the
project exists, in the language of the business or domain, not of the
software. It should be readable on its own.}

\subsection{Context and Overall Objective}\label{subsec:context-and-overall-objective}
\plt{One or two paragraphs: the domain the system lives in, the problem it
addresses, and the overall objective in a sentence a sponsor would recognize.}

\subsection{Current Situation}\label{subsec:current-situation}
\plt{How things are done today without the system, and what is wrong or
missing. Reference existing tools or processes the system will replace or
complement.}

\subsection{Expected Benefits}\label{subsec:expected-benefits}
\plt{The improvements the system will bring over the current situation, from
the point of view of the stakeholders. Where possible, make each benefit
measurable so it can be checked later.}

\begin{itemize}
    \item \textbf{Benefit} -- Description.
\end{itemize}

\subsection{Functionality Overview}\label{subsec:functionality-overview}
\plt{A high-level list of what the system does. Each item should correspond to
one or more components in S.1 and requirements in S.2, but stay at the level a
non-technical reader can follow.}

\begin{itemize}
    \item \textbf{Capability} -- Description.
\end{itemize}

\subsection{High-Level Usage Scenarios}\label{subsec:high-level-usage-scenarios}
\plt{The main ways users will interact with the system, as short narratives
(two or three sentences each). Detailed use cases belong in S.6. An activity
diagram can be included here if it helps.}

\begin{itemize}
    \item \textbf{Scenario 1: Name} -- Description.
\end{itemize}

% \begin{figure}[H]
%     \centering
%     \includegraphics[width=\textwidth]{activity_diagram.png}
%     \caption{Activity diagram showing the high-level usage scenarios.}
%     \label{fig:activity-diagram}
% \end{figure}

\subsection{Limitations and Exclusions}\label{subsec:limitations-and-exclusions}
\plt{What the system will not do. State exclusions explicitly so nobody later
assumes they were simply forgotten.}

\begin{itemize}
    \item \textbf{Exclusion} -- Description and rationale.
\end{itemize}

\subsection{Stakeholders and Requirements Sources}\label{subsec:stakeholders-and-requirements-sources}
\plt{Stakeholders: every person or group affected by the system, including
clients, users, supervisors, and the wider community. Requirements sources:
the documents, people, standards, and existing systems the requirements were
elicited from.}

\noindent\textbf{Stakeholders:}
\begin{itemize}
    \item \textbf{Stakeholder} -- Their interest in the system.
\end{itemize}

\noindent\textbf{Requirements Sources:}
\begin{itemize}
    \item \textbf{Source} -- What it contributed.
\end{itemize}

\subsection{User Profiles}\label{subsec:user-profiles}
\plt{For each distinct group of hands-on users, describe their domain
knowledge, goals, technical literacy, and which parts of the system matter most
to them. These profiles justify usability and accessibility requirements later
on.}

\noindent\textbf{1. Profile Name}
\begin{itemize}
    \item \textbf{Domain Knowledge:}
    \item \textbf{Goals:}
    \item \textbf{Technical Literacy:}
    \item \textbf{Rationale:} Which requirements exist for this user.
\end{itemize}

% ===========================================================================
\pegschapter{E}{Environment}

\plt{The Environment chapter describes the world the system operates in and
cannot change: physical constraints, other systems, rules of the domain. These
hold whether or not the system is built.}

\subsection{Glossary}\label{subsec:glossary}
The glossary has been moved to the front of the document (see the
\hyperref[sec:glossary]{Glossary} and the Changes from Original SRS Meyer
Template).

\subsection{Environment Components}\label{subsec:environment-components}
\plt{Elements of the environment that may affect or be affected by the system
and the project, including other systems the system must interface with,
hardware it runs on, and the people who operate it.}

\begin{itemize}
    \item \textbf{Component} -- Description and its relationship to the system.
\end{itemize}

\subsection{Constraints}\label{subsec:constraints}
\plt{Obligations and limits imposed on the system by the environment:
regulations, physical limits, resource limits, deadlines. These are not
design choices; the system must live with them.}

\begin{itemize}
    \item \textbf{Constraint} -- Description.
\end{itemize}

\subsection{Assumptions}\label{subsec:assumptions}
\plt{Properties of the environment that the system relies on but cannot
enforce. Each assumption is a risk if it turns out false; cross-reference the
Hazard Analysis where relevant.}

\begin{itemize}
    \item \textbf{Assumption} -- Description.
\end{itemize}

\subsection{Effects}\label{subsec:effects}
\plt{The ways the system changes its environment once deployed: effects on
users, on other systems, on resources.}

\begin{itemize}
    \item \textbf{Effect} -- Description.
\end{itemize}

\subsection{Invariants}\label{subsec:invariants}
\plt{Properties of the environment that must remain true throughout the
system's operation, and that the system must preserve.}

\begin{itemize}
    \item Invariant.
\end{itemize}

\subsection{Standards and Legal/Regulatory Factors}\label{subsec:standards}
\plt{Coding, documentation, interface or safety standards the project will
follow, and legal or regulatory factors such as licensing, intellectual
property, and data privacy. Cite each standard.}

\noindent\textbf{Standards:}
\begin{itemize}
    \item \textbf{Standard} -- What it governs and how it applies.
\end{itemize}

\noindent\textbf{Legal and Regulatory Factors:}
\begin{itemize}
    \item \textbf{Factor} -- Description.
\end{itemize}

\subsection*{Environmental Requirements}
\addcontentsline{toc}{subsection}{Environmental Requirements}

\plt{Requirements that follow directly from this chapter: installability on
the target platforms, resource limits, invariant enforcement, standards
compliance, legal compliance. Label each one so it can be referenced from S.7
and S.8.}

\begin{table}[H]
    \centering
    \caption{Environmental Requirements}
    \label{tab:env-requirements}
    \begin{tabular}{|l|p{11cm}|}
    \hline
    \textbf{Label} & \textbf{Requirement} \\
    \hline
    \label{NFR.E.1}NFR.E.1 & The system shall \ldots \\
    \hline
    \label{NFR.E.2}NFR.E.2 & The system shall \ldots \\
    \hline
    \end{tabular}
\end{table}

% ===========================================================================
\pegschapter{S}{System}

\plt{The System chapter is the technical core: what the system does, split
by component, with each requirement labelled, prioritized and verifiable.}

\subsection{Components}\label{subsec:components}
\plt{The major parts of the system and how they interact. Include a component
diagram. The components named here are reused as the subsections of S.2 and
the groupings of S.7.}

% \begin{figure}[H]
%     \centering
%     \includegraphics[width=0.8\textwidth]{component_diagram.png}
%     \caption{Component diagram.}
%     \label{fig:component-diagram}
% \end{figure}

\begin{itemize}
    \item \textbf{Component} -- Responsibility.
\end{itemize}

\subsection{Functionality}\label{subsec:functionality}
\plt{One subsubsection per component from S.1, each listing its functional
requirements with labels FR.S.<component>.<n>. Each requirement is a single,
testable ``shall'' statement. Non-functional requirements that cut across
components go in the last subsubsection as NFR.S.<n>.}

\subsubsection{Component Name}
\textbf{Functional Requirements}
\begin{itemize}
    \item \label{FR.S.1.1}\textbf{FR.S.1.1} The system shall \ldots
    \item \label{FR.S.1.2}\textbf{FR.S.1.2} The system shall \ldots
\end{itemize}

\subsubsection{Non-Functional Requirements}
\begin{itemize}
    \item \label{NFR.S.1}\textbf{NFR.S.1} \textbf{Quality attribute:} The system shall \ldots
\end{itemize}

\subsection{Interfaces}\label{subsec:interfaces}
\plt{The interfaces the system exposes or consumes: user interfaces, APIs,
protocols, file formats, hardware interfaces. For each, say who or what is on
the other side.}

\begin{itemize}
    \item \textbf{Interface} -- Description.
\end{itemize}

\subsection{Normal Operation}\label{subsec:normal-operation}
\plt{The expected flow of the system when everything goes right, from start-up
to shut-down. Keep it to the main path; alternatives and failures go in S.5
and S.6.}

\begin{enumerate}
    \item Step.
\end{enumerate}

\subsection{Undesired Event Handling}\label{subsec:undesired-event-handling}
\plt{How the system responds to failures, invalid inputs, and unexpected
conditions. Each entry should name the event and the required response.
Cross-reference the Hazard Analysis.}

\begin{itemize}
    \item \textbf{Event} -- Required response.
\end{itemize}

\subsection{Detailed Usage Scenarios}\label{subsec:detailed-usage-scenarios}
\plt{Full use cases for the scenarios introduced in G.5, one per scenario,
in the structure below. The main success scenario should be numbered so
extensions can reference its steps.}

\noindent\textbf{Use Case Name:} Name

\noindent\textbf{Primary Actor:} Actor

\noindent\textbf{Stakeholders and Interests:}
\begin{itemize}
    \item Stakeholder: interest.
\end{itemize}

\noindent\textbf{Preconditions:}
\begin{itemize}
    \item Precondition.
\end{itemize}

\noindent\textbf{Postconditions (Success Guarantees):}
\begin{itemize}
    \item Postcondition.
\end{itemize}

\noindent\textbf{Postconditions (Failure Guarantees):}
\begin{itemize}
    \item Postcondition.
\end{itemize}

\noindent\textbf{Main Success Scenario (Basic Flow):}
\begin{enumerate}
    \item Step.
\end{enumerate}

\noindent\textbf{Extensions:}
\begin{itemize}
    \item \textbf{Na.} Alternative or failure at step N, and how it is handled.
\end{itemize}

\noindent\textbf{Justification:} Why this scenario matters.

\subsection{Prioritization}\label{subsec:prioritization}
\plt{Every requirement from S.2 and the Environmental Requirements table,
with a MoSCoW priority and the reason for it. Group the tables by component.
The Date column records when the priority was last decided.}

\noindent Priorities follow MoSCoW: \textbf{M} must have, \textbf{S} should
have, \textbf{C} could have, \textbf{W} won't have (this release).

\subsubsection*{Component Name}
\begin{table}[H]
\centering
\begin{tabular}{|l|p{3.5cm}|c|p{5cm}|l|}
\hline
\textbf{ID} & \textbf{Requirement Name} & \textbf{Priority} & \textbf{Reasoning} & \textbf{Date} \\
\hline
\hyperref[FR.S.1.1]{FR.S.1.1} & Name & M & Reasoning & YYYY-MM-DD \\
\hline
\hyperref[FR.S.1.2]{FR.S.1.2} & Name & S & Reasoning & YYYY-MM-DD \\
\hline
\end{tabular}
\caption{Prioritization of Component Name requirements}
\label{tab:prioritization-component}
\end{table}

\subsection{Verification and Acceptance Criteria}\label{subsec:verification-and-acceptance-criteria}
\plt{How each class of requirement will be shown to be satisfied: unit and
integration tests, measurements, user studies, inspections. Make the criteria
concrete enough that the V\&V Plan can be derived from them. Every requirement
in S.2 should map to at least one criterion here.}

\noindent\textbf{Validation Overview:}

\noindent\textbf{Verification Approach by Requirement Class:}
\begin{itemize}
    \item \textbf{Class of requirements} -- How it will be verified and what counts as acceptance.
\end{itemize}

\subsection{Formalizations}\label{subsec:formalizations}
\plt{Mathematical or formal specification of the parts of the system that
benefit from precision: state models, invariants, algorithms, data formats.
Z notation is available through the \texttt{zed-csp} package. If nothing in
the system warrants formalization, say so and why.}

% ===========================================================================
\pegschapter{P}{Project}

\plt{The Project chapter covers how the system will be built: people,
imposed choices, schedule, deliverables, risks, and how requirements will be
managed over the project's life.}

\subsection{Roles and Personnel}\label{subsec:roles}
\plt{Each team member and their role, plus the supervisor and TA. Keep this
consistent with the Development Plan.}

\begin{itemize}
    \item \textbf{Role (Name)} -- Responsibilities.
\end{itemize}

\subsection{Imposed Technical Choices}\label{subsec:technical-choices}
\plt{Technologies that are fixed before design begins and why: languages,
frameworks, hardware, platforms mandated by the client, supervisor, or
environment. Choices the team is still free to make do not belong here.}

\begin{itemize}
    \item \textbf{Choice} -- What is imposed and by whom.
\end{itemize}

\subsection{Schedule and Milestones}\label{subsec:schedule}
\plt{The major milestones and their dates. Course deliverables are a starting
point; add project-specific milestones such as prototypes or integration
points.}

\begin{table}[H]
  \centering
  \begin{tabular}{|l|l|}
  \hline
  \textbf{Milestone} & \textbf{Due Date} \\ \hline
  Milestone & Date \\ \hline
  \end{tabular}
  \caption{Project Milestones and Due Dates}
  \label{tab:project-milestones}
\end{table}

\subsection{Tasks and Deliverables}\label{subsec:tasks}
\plt{The work breakdown: each task, what it produces, and which milestone it
feeds. Tasks should be small enough to assign to one or two people.}

\begin{table}[H]
\centering
\begin{tabular}{|p{3cm}|p{5.5cm}|p{3cm}|p{2.5cm}|}
\hline
\textbf{Task} & \textbf{Description / Steps} & \textbf{Deliverable} & \textbf{Due Date} \\ \hline
Task & Description & Deliverable & Date \\ \hline
\end{tabular}
\caption{Tasks and Deliverables}
\label{tab:tasks-deliverables}
\end{table}

\subsection{Required Technology Elements}\label{subsec:required-tech}
\plt{Everything the project needs to obtain or build before the system can be
delivered: libraries, hardware, datasets, services, licences, compute. Note
anything not yet secured.}

\begin{itemize}
    \item \textbf{Element} -- Purpose and status.
\end{itemize}

\subsection{Risk and Mitigation Analysis}\label{subsec:risk}
\plt{Project risks (schedule, technical, people, external dependencies), each
with its likelihood, impact, and mitigation. System hazards belong in the
Hazard Analysis, not here.}

\begin{itemize}
    \item \textbf{Risk: Name}
    \begin{itemize}
        \item Description:
        \item Mitigation:
    \end{itemize}
\end{itemize}

\subsection{Requirements Process and Report}\label{subsec:requirements-process}
\plt{How requirements were elicited and how they will be maintained: who was
consulted, through which channels, how changes are tracked, and how this
document will be kept in sync with the Hazard Analysis, V\&V Plan and design.}

\begin{itemize}
    \item \textbf{Elicitation} --
    \item \textbf{Communication} --
    \item \textbf{Change Tracking} --
    \item \textbf{Process Updates} --
\end{itemize}

% ===========================================================================
\newpage
\tocsection{Safety and Security Requirements}
\label{sec:safety-security}

\plt{Port the requirements from the Safety and Security Requirements section
of the Hazard Analysis so this document is the single list of requirements.
Keep the labels identical in both documents.}

\subsection*{Functional Safety Requirements}
\begin{itemize}
    \item \label{FR.Sa.1}\textbf{FR.Sa.1} The system shall \ldots
\end{itemize}

\subsection*{Non-Functional Safety Requirements}
\begin{itemize}
    \item \label{NFR.Sa.1}\textbf{NFR.Sa.1} The system shall \ldots
\end{itemize}

% ===========================================================================
\newpage
\tocsection{Requirement Dependencies}
\label{sec:requirement-dependencies}

\plt{A traceability diagram or matrix showing which requirements depend on
which. Use the landscape environment below for a wide figure.}

% \begin{landscape}
%     \centering
%     \includegraphics[width=1.5\textheight]{requirement_dependency_diagram.png}
%     \captionof{figure}{Requirement dependency diagram.}
%     \label{fig:requirement-dependencies}
% \end{landscape}

% ===========================================================================
\newpage
\plt{Cite sources with \texttt{\textbackslash citep\{key\}} or
\texttt{\textbackslash citet\{key\}}; entries live in
\texttt{refs/References.bib}. The References section appears once at least
one source is cited.}

\bibliographystyle{plainnat}
\bibliography{../../refs/References}

% ===========================================================================
% Appendices use ordinary section numbering again.
\renewcommand{\thesubsection}{\thesection.\arabic{subsection}}
\renewcommand{\theHsubsection}{\thesection.\arabic{subsection}}
\newpage
\section{Appendix --- Generative AI Use Disclosure}

\wss{ChatGPT (version 5) was used to brainstorm ideas for the template for this disclose section.}

\subsection{AI Tools Used}

\wss{Give the name of the AI tool(s) used and the version.}

\subsection{How and Where Used}

\wss{Explain what the tool(s) were used for.  Examples include: brainstorming,
explaining a technical concept, reviewing the document, generating or modifying
text, generating or modifying diagrams, identifying errors or omissions or
inconsistencies.}

\wss{For each of the uses, identify which parts of the document were affected.}

\wss{You do not need to report every interaction, but you should disclose the
uses that materially contributed to your work}

\subsection{Verification of AI-Generated Content}

\wss{\begin{itemize}
    \item Describe how the team checked AI-generated or AI-assisted material
    before including it in the document.
    \item In particular, verify factual claims, technical assertions,
    requirements, calculations, references, and other information for which an
    AI system may produce incorrect or fabricated results.
    \item \textbf{The team is responsible for the accuracy, completeness,
consistency, and appropriateness of the submitted document, regardless of
whether material was generated by a human or an AI system.}
\end{itemize}}

\newpage

\section{Appendix --- Reflection}

\input{../Reflection.text}

\input{../SRS_Reflection.text}

\end{document}
